\documentclass[a4paper,11pt]{amsart}
\usepackage[T1]{fontenc}
\usepackage{lmodern,amsmath,amssymb,amsthm,mathtools}
\usepackage[margin=28mm]{geometry}
\usepackage{microtype,booktabs,array,graphicx}
\usepackage[hidelinks,pdfusetitle]{hyperref}
\newcommand{\repofile}[2]{\href{https://github.com/MColbrook/transmission-shape-optimisation/blob/main/#1}{#2}}
\numberwithin{equation}{section}
\theoremstyle{plain}
\newtheorem{theorem}{Theorem}[section]
\newtheorem{proposition}[theorem]{Proposition}
\newtheorem{lemma}[theorem]{Lemma}
\newtheorem{corollary}[theorem]{Corollary}
\theoremstyle{definition}
\newtheorem{remark}[theorem]{Remark}
\newcommand{\R}{\mathbb R}
\newcommand{\C}{\mathbb C}
\newcommand{\D}{\mathbb D}
\newcommand{\Hcl}{H_0^2}
\newcommand{\dd}{\,\mathrm d}
\newcommand{\trspec}{\sigma_{\mathrm{tr}}}
\DeclareMathOperator{\Rea}{Re}
\DeclareMathOperator{\Ima}{Im}
\DeclareMathOperator{\supp}{supp}
\DeclareMathOperator{\cof}{cof}
\DeclareMathOperator{\diag}{diag}
\DeclareMathOperator{\dist}{dist}

\title[Mathematical supplement]{Mathematical supplement to\\
Shape optimisation of the least-modulus transmission eigenvalue}
\author{Matthew J. Colbrook}
\author{Catherine Drysdale}
\hypersetup{pdfauthor={Matthew J. Colbrook and Catherine Drysdale}}
\begin{document}
\raggedbottom
\maketitle

Can the same domain beat the disk throughout an interval of
contrasts? In Section~\ref{sup:trial} we construct a smooth
domain of area \(\pi\) with a transmission eigenvalue of smaller
modulus than every disk eigenvalue for each \(n\in[5.18,5.281738]\).
Section~\ref{sup:hessian} gives the Fourier coefficient bounds for
the disk Hessian below resonance, including at unit contrast.
Above resonance, the comparison for one-signed minimisers depends
on the signs we prove in Section~\ref{sup:radial}.

Let \(\D=\{z\in\mathbb C:|z|<1\}\). We denote the \(k\)-th positive
zero of the Bessel function \(J_m\) by \(j_{m,k}\) and set
\[
 \alpha=j_{0,1}^2,\qquad \beta=j_{1,1}^2,\qquad
 \zeta=j_{0,2}^2,\qquad N=\zeta/\alpha.
\]
For a real field \(u\) on \(\Omega\), let \(I_0,I_1,I_2\) be the squared
\(L^2(\Omega)\) norms of \(u,\nabla u,\Delta u\), respectively.
The form at the spectral parameter \(t\) is
\(q_{\Omega,n}(t;u)=I_2-(n+1)tI_1+nt^2I_0\).

The \href{https://github.com/MColbrook/transmission-shape-optimisation}{companion repository}
contains the exact data and code. The \repofile{checker-map.md}{certificate index}
identifies the calculation for each inequality, and the
\repofile{README.md}{README} gives the verification commands.

\section{Exact trial data}\label{sup:trial}

We start with a polynomial image of the disk. Rational coefficients
allow us to bound the trial form by exact arithmetic after we
normalise the area. Define
\begin{equation}\label{eq:map12}
 F(z)=z+10^{-8}\sum_{j=0}^{11}a_jz^{2j+3},
 \qquad f(z)=S^{-1/2}F(z),\qquad \Omega_{12}=f(\D),
\end{equation}
where the integer coefficients are
\[
\begin{split}
(a_j)={}&(23661008,\ 944390,\ 841714,\ -267345,\ 182938,\ -115928,\\
        &\hspace{10mm}80287,\ -57758,\ 42974,\ -32778,\ 25473,\ -19906)
\end{split}
\]
and
\[
 S=1+10^{-16}\sum_{j=0}^{11}(2j+3)a_j^2
   =\frac{11690370621775763}{10^{16}}.
\]
Let $g=F'$. The coefficient sum gives
\[
 |g(z)-1|\le\eta:=\frac{92297023}{10^8}<1
 \quad (|z|\le1).
\]
Hence \(\Rea g>0\) on a neighbourhood of the closed disk.
Integrating \(g\) along a chord shows that \(F\) is injective.
Since \(g\) does not vanish, the image boundary is analytic.
The conformal area formula gives \(|\Omega_{12}|=\pi\).

For a real clamped trial \(U\) on \(\D\), let \(u=U\circ f^{-1}\).
We compute its moments \(I_j=I_j(u)\) on \(\Omega_{12}\) by changing
variables back to the disk.
\begin{equation}\label{eq:conformal-forms}
 I_0=\frac1S\int_\D U^2|g|^2,\qquad
 I_1=-\int_\D U\Delta U,\qquad
 I_2=S\int_\D\frac{(\Delta U)^2}{|g|^2}.
\end{equation}
Set \(s=r^2\) and take the trial
\begin{equation}\label{eq:trial12}
 U=(1-s)^2\sum_{j=0}^{6}\sum_{k=0}^{4}
       \frac{b_{jk}}{10^{10}}\,
       r^{2j+1}\cos((2j+1)\theta)
       P_k^{(4,2j+1)}(2s-1),
\end{equation}
with the integers \(b_{jk}\) in Table~\ref{tab:trial}.
The factor \((1-s)^2\) makes both the trace and the normal
derivative vanish on the boundary. For the Jacobi polynomials in
\eqref{eq:trial12}, we use
\[
 P_k^{(4,m)}(2s-1)
 =\sum_{i=0}^k
 \binom{k+4}{i}\binom{k+m}{k-i}s^i(s-1)^{k-i}.
\]

\begin{table}[!ht]
\centering\small
\caption{Integer numerators \(b_{jk}\) for the clamped trial
\eqref{eq:trial12} on the unit disk. The Jacobi basis coefficients
are \(b_{jk}/10^{10}\).}\label{tab:trial}
\begin{tabular}{r r r r r r}
\toprule
$j$&$b_{j0}$&$b_{j1}$&$b_{j2}$&$b_{j3}$&$b_{j4}$\\
\midrule
0&10000000000&-714354563&-56400458&-2125776&1583967\\
1&3570773036&-529889309&23843272&-597204&292049\\
2&785588028&-133167479&9457497&-308100&59819\\
3&134049214&-29915606&2893523&-135368&8429\\
4&29725599&-4266203&416310&-17814&2314\\
5&-1030520&-1130923&153893&-11423&128\\
6&6047546&300537&6035&2833&353\\
\bottomrule
\end{tabular}
\end{table}

We integrate polynomials using the exact disk moments
\begin{equation}\label{eq:moments}
 \frac1\pi\int_\D z^p\overline z^{\,q}
 =\frac{\mathbf1_{\{p=q\}}}{p+1}.
\end{equation}
The reciprocal in \(I_2\) is the only term we approximate. We use
the recurrence for \(g^{-1}\) to compute an even polynomial \(H\)
of degree at most \(600\), rounding each step to a multiple of
\(10^{-24}\). We then compute the residual \(R=1-gH\) exactly
and set
\[
 R_*=\sum_k|[z^k]R|,\qquad e=\frac{R_*}{1-\eta}.
\]
Then
\[
 \|g^{-1}-H\|_{L^\infty(\D)}\le e,\qquad
 \frac{I_2}{\pi}\le A_+:=
 \frac S\pi\int_\D|\Delta U\,H|^2+
 \frac S\pi\int_\D(\Delta U)^2
       \left(\frac{2e}{1-\eta}+e^2\right).
\]
The moment formula \eqref{eq:moments} and the residual bound give
\begin{equation}\label{eq:compact-cert}
 \frac{I_1}{\pi}>\frac{2941129141}{10^{10}},
 \qquad
 \frac{I_0}{\pi}<\frac{7471409}{400000000},
 \qquad
 A_+<\frac{6923409}{10^6},
 \qquad e<10^{-10}.
\end{equation}
The \repofile{certificates/shape_optimization/certify_near_optimal_domain.py}{trial calculation}
uses rational arithmetic in the Python standard library.

The saved trial certificate checks the form at \(t=51963/10000\)
and \(n=263447/50000\). For the comparison at resonance, we use
the same \(t\) and the slightly smaller contrast
\(n_0=526893/100000\). The same moment bounds
\eqref{eq:compact-cert} give
\[
 \frac{q_{\Omega_{12},n_0}(t;u)}{\pi}
 < -\frac{47050322517501147}{4\cdot10^{21}}.
\]
At this trial and value of \(t\), we have \(t^2I_0-tI_1<0\) by
\eqref{eq:compact-cert}, so the form decreases with \(n\).
The endpoint calculation in the next subsection gives the Bessel
bounds \(\alpha<5.78319\) and \(\zeta>30.47126\). Thus
\(N>30.47126/5.78319>n_0\) and the form at \(t=5.1963\) is negative
at resonance. The crossing lemma in the paper then gives a positive
real eigenvalue below \(5.1963\).

\subsection{The contrast sweep and endpoint}
How far from resonance does the same domain improve on the disk?
We first look for a bound that lies above the first positive real
eigenvalue of \(\Omega_{12}\) and below the modulus of every disk
eigenvalue. On the closed contrast interval
\([259/50,5281/1000]\), set
\[
 T(n)=\frac{519631}{100000}
      -\frac54\left(n-\frac{263447}{50000}\right)
      +\frac12\left(n-\frac{263447}{50000}\right)^2 .
\]
The moment bounds \eqref{eq:compact-cert} give a degree-five upper
bound for \(q_{\Omega_{12},n}(T(n);u)/\pi\). Each Bernstein
coefficient of this polynomial on the interval is less than
\(-2.64\cdot10^{-6}\), so the form is negative throughout.

For the disk we start with the radial sector. We truncate the
entire function \(E_n\) from Section~\ref{sup:hessian} at degree
\(32\). For \(1\le n\le6\) and \(|z|\le8\), the term of degree \(k\)
has modulus at most \(28^k/(k!)^2\). Its first partial derivatives
are bounded by \(k\) times this quantity. We bound each tail by
a geometric series beginning at degree \(33\).

The interval calculation encloses
\(E_n(T(n)(1-s^2+2is)/(1+s^2))\) on \(1561\) rectangles in the
contrast and \(s\) variables, with \(0\le s\le1\). Every enclosure
excludes zero. This gives the zero-free contour used for the radial
count in the paper. The count is zero at resonance and hence remains
zero inside \(|z|=T(n)\) throughout the interval.

For nonreal values, the angular lower bound excludes every sector
of order \(m\ge1\). For positive real values, we check the first
two angular determinants on \(29\) contrast cells and use the Bessel
signs to exclude the higher sectors. The
\repofile{certificates/continuation/certify_contrast_interval.json}{interval data}
give the exact partitions and the bounds for both interval tests.

At the upper endpoint \(n_b=5.281738\), we test the trial at the
active radial disk eigenvalue \(\ell(n_b)\). In the
\repofile{certificates/check_transition_endpoints.mjs}{endpoint calculation},
we bracket the ordinary Bessel series by the alternating partial
sums through indices \(30\) and \(31\). For the modified series,
the squared Bessel argument lies in \((0,31)\), and the tail beyond
index \(30\) is at most \(100/99\) times the first omitted term.
We obtain
\[
 5.1803707<\ell(5.281738)<5.1803709,\qquad
 t_D(5.281738)<142,
\]
where \(t_D(n)=n\ell(n)^2\) for the active radial eigenvalue
\(\ell(n)\) above resonance. This bracket and the trial moment
bounds give an upper bound less than \(-3.28\cdot 10^{-6}\) for
\(q_{\Omega_{12},n_b}(\ell(n_b);u)/\pi\). The slope of the disk
envelope carries the strict comparison back to resonance.
Together with the contour calculation, this covers the closed
interval \([5.18,5.281738]\).

The same series bounds give \(10.21582622<\gamma<10.21582624\)
for the square root of the first clamped-plate value on \(\D\).
Comparing \(t_D(n)\) with \(\gamma^2\) brackets the plate cutoff in
\((5.4254619,5.4254621)\). The integer inequality
\(104\cdot207>142^2\) is the product estimate used in the paper's
plate gap argument.

\section{Subresonant Hessian certificates}\label{sup:hessian}

Let \(\lambda_0\) be the upper member of the least-modulus disk
pair for \(1\le n<N\). The disk count in the paper gives at most one
eigenvalue in the upper half-plane with
\(\lvert\sqrt n\,\lambda\rvert\le14\). Thus, to identify \(\lambda_0\)
on a contrast interval, it suffices to enclose a radial zero in this
region, where the paper places the upper member of the least-modulus
pair. At \(n=1\) we use the clamped pencil \((\Delta+\lambda)^2\).

Set \(x=\lambda_0\), \(y=n\lambda_0\) and
\(\kappa=\sqrt{x}J_0'(\sqrt{x})/J_0(\sqrt{x})\).
For \(m\ge1\), let \(V_m\) be the radial factor of order \(m\) with
\(V_m(1)=0\) and \(V_m'(1)=1\). For \(n>1\), we solve for the
coefficients of \(J_m(\sqrt{x}\,r)\) and \(J_m(\sqrt{y}\,r)\) from
these two conditions. The determinant of this system is nonzero
since the kernel at \(\lambda_0\) is radial. For
\(B_m=V_m''(1)+1-\kappa\), this gives
\begin{equation}\label{um:B-def}
 B_m=-\kappa
 -\frac{(y-x)J_m(\sqrt x)J_m(\sqrt y)}
 {\sqrt yJ_m'(\sqrt y)J_m(\sqrt x)
       -\sqrt xJ_m'(\sqrt x)J_m(\sqrt y)}.
\end{equation}
At \(n=1\), we take the confluent value of \eqref{um:B-def} using
the basis obtained by a divided difference of the radial solutions
at \(x\) and \(y\).

To enclose the radial zeros, we use the normalised determinant
\[
 E_n(\lambda)
 =\sum_{k=0}^\infty\frac{(-\lambda/4)^k}{k+1}
             \sum_{j=0}^k\frac{n^j}{j!^2(k-j)!^2}.
\]
For \(n>1\), the radial Wronskian identity gives the boundary
determinant \(D_0=(1-n)\lambda E_n/2\). Integrating the Bessel
products term by term gives the series. At \(n=1\), dividing
the Wronskian by \(1-n\) before taking the limit gives the boundary
determinant of the confluent radial basis.

Set \(t=\sqrt n\), \(\mu=t\lambda\) and \(u=t+t^{-1}\). We
expand the symmetric Laurent coefficients of \(E_{t^2}(\mu/t)\)
in powers of \(u\) to obtain the degree-\(24\) truncation
\begin{equation}\label{uh:Q}
 Q(\mu,u)=\sum_{k=0}^{24}\sum_{j=0}^{\lfloor k/2\rfloor}
       \frac{(-\mu)^k u^{k-2j}}
              {4^k(k+1)k!(j!)^2(k-2j)!}.
\end{equation}
To bound the error for \(t|\mu|\le b\), we use the unit-contrast
series with positive coefficients. Its coefficient of \(b^k\) is
\(\binom{2k}{k}/[4^k(k+1)(k!)^2]\). For \(d\le3\), the
order-\(d\) derivative tail at \(b\), multiplied by \(t^d\),
bounds the error after \(d\) differentiations in \(\mu\). In that
tail, the ratio of the terms of original indices \(k+1\) and \(k\)
is at most \(b/[(k+2)(k+1-d)]\) for \(k\ge25\). This allows us to
bound each tail by a geometric series.

The numerator and denominator in \eqref{um:B-def} both vanish
as \(y\to x\). We use an integral formula that also applies at
unit contrast. Let
\[
 T_m(z)=\sum_{j=0}^\infty\frac{(-z/4)^j m!}{j!(m+j)!},
 \qquad
 L_m=2\int_0^1r^{2m+1}T_m(xr^2)T_m(yr^2)\,\mathrm dr .
\]
Thus \(T_m(z)=m!(2/\sqrt z)^m J_m(\sqrt z)\), with the value at
zero given by the series. The Wronskian identity now gives
\begin{equation}\label{um:B-integral}
 B_m=-\kappa+\frac{2T_m(x)T_m(y)}{L_m},
 \qquad
 \kappa=-\frac{x T_1(x)}{2T_0(x)}.
\end{equation}
The term of order \(k\) in \(L_m\) equals
\[
 \frac{(-\mu/4)^k(m!)^2}{m+k+1}
   \sum_{j=0}^k\frac{t^{2j-k}}
        {j!(k-j)!(m+j)!(m+k-j)!}.
\]
Pairing reciprocal powers of \(t\) shows that the absolute
coefficient is nondecreasing for \(t\ge1\). The radial coefficient
bounds it from above. For \(T_m\), the ratio of successive terms
at \(|z|\le b\) is bounded by \(b/[4(k+1)(m+k+1)]\) for \(k\ge31\).
We use this ratio for the tail of its degree-\(30\) truncation.

We will use the Bessel product to bound the modes of high order.
Matching the first two Taylor coefficients gives, for \(A_j=j_{m,j}^2\),
\begin{equation}\label{uh:rayleigh}
 \sum_j A_j^{-1}=\frac{1}{4(m+1)},\qquad
 \sum_j A_j^{-2}=\frac{1}{16(m+1)^2(m+2)}.
\end{equation}
In particular, \(A_1>4(m+1)\sqrt{m+2}\).
The first angular order comes from translations. Differentiating
the radial state in a Cartesian direction gives the field with
radial factor \(V_1\) once we normalise its derivative trace. With
the radial mode \(U\) normalised by \(U''(1)=1\), we have
\(U'''(1)=\kappa-1\), so \(B_1=0\), including at unit contrast.

\subsection{The first compact interval}
\begin{proposition}[Extended all-mode Hessian bound]
\label{ux:hessian}
For every \(1\le n\le441/100\), one has \(B_1(n)=0\) and
\[
 \operatorname{Re}B_m(n)>\frac14(m+1),\qquad
 |B_m(n)|<14(m+1),\qquad m\ge2.
\]
\end{proposition}

\begin{proof}
We use the \(550\) closed rational intervals in
\path{subresonant/unit_midrange/extension/centres.json}. These
cover \([1,21/10]\) in \(t\) and adjoining intervals share an
endpoint. On each interval we enclose the root in a circle with
rational centre \(z_c\) chosen numerically. The interval lengths
and radii are
\[
\begin{array}{c|c|c}
 \text{parameter range}&\text{maximum length}&\text{radius }r\\
 {[1,3/2]}&1/200&1/50\\
 {[3/2,19/10]}&1/500&3/200\\
 {[19/10,41/20]}&1/1000&1/100\\
 {[41/20,21/10]}&1/2000&3/400.
\end{array}
\]
We bound the residual by \(\epsilon\) using the value of
\(Q(z_c,u)\) at the midpoint in \(u\), its \(u\)-derivative on the
whole interval and the positive series tail at argument \(30\).
We evaluate \(\partial_\mu Q(z_c,u)\) over the same interval and
subtract the derivative tail from the resulting lower bound. This
gives a positive lower bound \(d_1\) for the modulus of the
derivative of the full series. To bound the second derivative,
we take a rational upper bound \(s_c\) for \(|z_c|\) and use a
radius \(R_c\) satisfying
\[
 R_c=s_c+2r<14,\qquad t_+R_c<25.
\]
The absolute coefficient sum at \(R_c\), together with the second
derivative tail, gives an upper bound \(d_2\) throughout the root
circle. In these bounds the order-\(j\) tails carry the factor
\((21/10)^j\). On every interval we verify
\[
 \epsilon+\frac12r^2d_2<rd_1.
\]
Taylor's formula and Rouch\'e's theorem therefore give one zero
in \(|\mu-z_c|<r\) for every parameter in the interval. The checks
\(\operatorname{Im}z_c>r\) and \(R_c<14\) place that zero in the
region of the disk count, which identifies the upper disk value.

We now enclose \(\mu\) in the square
\[
 |\operatorname{Re}(\mu-z_c)|\le r,\qquad
 |\operatorname{Im}(\mu-z_c)|\le r.
\]
We still have \(|\mu|<R_c\) on the square. Using the intervals
\([1/t_+,1/t_-]\) and \([t_-,t_+]\) gives enclosures for
\(x=\mu/t\) and \(y=t\mu\), with
\[
 |x|,|y|<25.
\]
In \eqref{um:B-integral}, we evaluate \(T_m\) to degree \(30\)
with a geometric tail at argument \(25\). For \(L_m\) we use
degree \(24\) and the radial tail at argument \(30\).
The coefficient magnitudes of \(L_m\) are nondecreasing in
\(t\ge1\), so the endpoint values bound them on each interval.
With rational complex intervals we obtain
\[
 \operatorname{Re}\kappa_n<3,\qquad |\kappa_n|<10,
\]
and the asserted inequalities for \(2\le m\le6\). Each
denominator rectangle excludes zero. The interval coverage and
all these comparisons are verified in
\path{subresonant/unit_midrange/extension/check_hessian.py}.

For \(m\ge7\), let \(A_j=j_{m,j}^2\). The Rayleigh sum \eqref{uh:rayleigh} gives
\[
 A_j\ge A_1>4(m+1)\sqrt{m+2}\ge12(m+1)\ge96.
\]
For
\[
 S_m=\sum_j\frac{A_j}{(A_j-x)(A_j-y)},\qquad
 s_m=\frac1{4(m+1)},
\]
the Bessel product gives \(B_m=-\kappa_n+1/(2S_m)\).
Since \(|x|,|y|<25\),
\[
 |S_m-s_m|\le d s_m,\qquad
 d=\left(1-\frac{25}{96}\right)^{-2}-1
   =\frac{4175}{5041}<1.
\]
Inverting this disk proves
\[
 \operatorname{Re}B_m
 \ge\frac{2(m+1)}{1+d}-3>\frac14(m+1),
 \qquad
 |B_m|\le10+\frac{2(m+1)}{1-d}<14(m+1).
\]
The strict comparisons hold at \(m=7\) by rational arithmetic
and improve with \(m\). The translation identity gives \(B_1=0\).
\end{proof}

\subsection{A compact cone comparison}
\label{uh:compact-section}

On the next interval we bound the arguments of \(\kappa\)
and \(x/\kappa\). Use
\[
 x=\lambda_0(n),\qquad y=n\lambda_0(n),\qquad
 \mu=\sqrt n\,\lambda_0(n),\qquad
 \kappa=\frac{\sqrt xJ_0'(\sqrt x)}{J_0(\sqrt x)},\qquad
 Z=\frac{x}{\kappa}.
\]
Neither denominator vanishes, since the nonzero zeros of the
Bessel functions are real and \(x\) is nonreal.

\begin{proposition}[Compact cone bounds]
\label{uh:compact-cones}
For
\[
 \frac{441}{100}\le n\le
       \left(\frac{459083}{200000}\right)^2,
\]
the active radial pair satisfies
\[
 -\operatorname{Re}\kappa>\frac{|\kappa|}{8},\qquad
 -\operatorname{Re}Z>\frac{|Z|}{3},\qquad
 \operatorname{Re}(Z^2)<0.
\]
\end{proposition}

\begin{proof}
We use the \(808\) closed rational intervals in
\repofile{subresonant/unit_midrange/cap/hessian_compact_centres.json}{centre data}.
Adjoining intervals share an endpoint, and together they cover
\[
 21/10\le t=\sqrt n\le459083/200000,
\]
The centres \(z_c\) are chosen numerically. To enclose a zero
throughout each interval, we use the polynomial \(Q(\mu,u)\) in
\eqref{uh:Q}. The residual at \(z_c\) is bounded by
\[
 |Q(z_c,u_c)|
 +\frac{u_+-u_-}{2}
       \sup_{u_-\le u\le u_+}|\partial_uQ(z_c,u)|+\tau_0.
\]
Here \(u_c\) is the midpoint of \([u_-,u_+]\) and \(\tau_j\)
bounds the order-\(j\) derivative tail of the positive unit-contrast
series at argument \(31\). We use the same geometric estimate as
before, with a factor at most \((23/10)^j\) for differentiation in
\(\mu\). On every interval we have
\[
 |z_c|+2r<53/4,\qquad
 t_+(|z_c|+2r)<31,\qquad \operatorname{Im}z_c>r.
\]

To bound the modulus of the spectral derivative from below, we
take the larger of the lower bounds for the absolute real and
imaginary parts of \(\partial_\mu Q(z_c,u)\) and subtract the
derivative tail. This gives a positive bound \(d_1\). For the
second derivative, we use its value at the centre and a bound on
the third derivative throughout the circle. The resulting bound is
\[
 d_2=\sup_{u_-\le u\le u_+}
          |\partial_{\mu\mu}Q(z_c,u)|
       +(23/10)^2\tau_2+rM_3,
\]
where \(M_3\) is an absolute coefficient bound for the full
third derivative on \(|\mu|\le|z_c|+2r\).
With \(\epsilon\) the residual bound, we verify on every interval that
\[
 \epsilon+\frac12r^2d_2<rd_1.
\]
Taylor's formula and Rouch\'e's theorem give one upper
root in \(|\mu-z_c|<r\) for every parameter in the interval.
The disk count again identifies the upper disk value.

We enclose the root in the square of half-width \(r\) and use
the interval for \(1/t\) to get a rectangle for \(x\).
For the normalised Bessel functions
\[
 T_m(s)=m!\left(\frac2{\sqrt s}\right)^mJ_m(\sqrt s)
\]
we use degree \(30\) and the geometric tails at argument \(31\)
to evaluate
\[
 \kappa=-\frac{xT_1(x)}{2T_0(x)},\qquad
 Z=-\frac{2T_0(x)}{T_1(x)}.
\]
The second formula cancels the common factor \(x\) before
interval evaluation. Each denominator rectangle excludes zero.
With \(K\) and \(W\) the resulting rectangles, we verify
\[
 \sup\operatorname{Re}K<0,\qquad
 63\bigl(-\sup\operatorname{Re}K\bigr)^2
       >\sup|\operatorname{Im}K|^2,
\]
\[
 \sup\operatorname{Re}W<0,\qquad
 8\bigl(-\sup\operatorname{Re}W\bigr)^2
       >\sup|\operatorname{Im}W|^2,
 \qquad \sup\operatorname{Re}(W^2)<0.
\]
These inequalities give the three bounds in the proposition.
The rational comparisons, including the interval coverage, are
verified in \path{subresonant/unit_midrange/cap/hessian_compact_check.py}.
\end{proof}
\subsection{Cone bounds up to resonance}
\label{uhe:section}

The logarithmic derivative \(\kappa\) has a pole at \(x=\alpha\).
To obtain cone bounds up to resonance, we use \(x=\lambda_0(n)\)
and set
\[
 \kappa=\frac{\sqrt xJ_0'(\sqrt x)}{J_0(\sqrt x)},\qquad
 \eta=(N-n)^{1/3}.
\]
We enclose the root in a fixed disk in the variable
\(z=(x-\alpha)/\eta\). The resulting contrast interval overlaps
the interval in Proposition~\ref{uh:compact-cones}.

\begin{proposition}[Endpoint cone bounds]\label{uhe:cones}
For every \(0<\eta\le1/40\), with \(n=N-\eta^3\), one has
\[
 -\operatorname{Re}\kappa>\frac{|\kappa|}{6},\qquad
 -\operatorname{Re}\frac{x}{\kappa}>
                 \frac14\left|\frac{x}{\kappa}\right|,\qquad
 \operatorname{Re}\!\left[\left(\frac{x}{\kappa}\right)^2\right]<0.
\]
Moreover,
\[
 x=\alpha+\eta z,\qquad
 \left|z-\left(-1.30655+2.26302\,i\right)\right|<\frac1{10}
\]
throughout this interval. All stated decimal values are exact
terminating rational numbers.
\end{proposition}

\begin{proof}
We first enclose the squared Bessel zeros \(\alpha\) and \(\zeta\)
by checking the signs at the endpoints of the brackets
\[
 \begin{split}
 5.78318596294&<\alpha<5.78318596296,\\
 30.47126234365&<\zeta<30.47126234368.
 \end{split}
\]
We use degree-\(40\) Bessel series with tails bounded by the
first omitted absolute term times \((1-b/(4\cdot42^2))^{-1}\),
where \(b\le31\). The known zero ordering identifies \(\alpha\)
and \(\zeta\) as the squared zeros in these brackets.

To expand the determinant near resonance, we work with Bessel
functions normalised at each zero. For \(a=\alpha,\zeta\), let
\[
 f_a(h)=
 \frac{J_0(\sqrt{a+h})}
      {\left.\frac{\mathrm d}{\mathrm dz}J_0(\sqrt z)\right|_{z=a}}.
\]
The zero at \(a\) is simple, so \(f_a\) is well defined and
satisfies \(f_a(0)=0\) and \(f_a'(0)=1\). The Bessel equation gives
\[
 (a+h)f_a''+f_a'+\frac14f_a=0.
\]
Substituting \(f_a(h)=\sum_{j\ge0}c_j(a)h^j\) gives the recurrence
\[
 c_0=0,\quad c_1=1,\qquad
 a(j+2)(j+1)c_{j+2}
 =-(j+1)^2c_{j+1}-\frac14c_j.
 \tag{UHE1}\label{uhe:recurrence}
\]
In particular,
\[
 f_a(h)=h-\frac{h^2}{2a}
        +\left(\frac1{3a^2}-\frac1{24a}\right)h^3+O(h^4).
\]

Set \(x=\alpha+\eta z\), \(y=nx\), and
\[
 h_x=\eta z,\qquad
 h_y=N\eta z-\alpha\eta^3-\eta^4z.
\]
The radial characteristic equation is
\[
 D(\eta,z):=
 x f_\alpha'(h_x)f_\zeta(h_y)
 -y f_\zeta'(h_y)f_\alpha(h_x)=0.
\]
Using \(\zeta=N\alpha\), we find that the coefficients of
\(\eta^0,\eta^1,\eta^2\) vanish identically. Substituting the
coefficients from \eqref{uhe:recurrence} gives
\[
 \frac{D(\eta,z)}{\eta^3}
 =\frac{N(N-1)}{12}z^3-\alpha^2+O(\eta).
 \tag{UHE2}\label{uhe:normalised-D}
\]
Indeed, \(x f_\alpha'(h_x)=\alpha-h_x^2/8+O(h_x^3)\).
The coefficient of \(\eta^3z^3\) in \(D\) is
\[
 \alpha N^3c_3(\zeta)-\zeta c_3(\alpha)
       +\frac{N^2-N}{8}
 =\frac{N(N-1)}{12},
\]
while the constant part at that order is \(-\alpha^2\).
Thus the quotient extends analytically to \(\eta=0\). We use
the cancellation of the first three coefficients before evaluating
the quotient in intervals.

To bound the determinant of the truncated series, we first need
uniform bounds on the tails of \(f_a\) and \(f_a'\). Let \(a_-\)
be the lower endpoint of the bracket for \(a\), and let
\[
 L_a=\max\{|c_{11}(a)|a_-^{10},\,|c_{12}(a)|a_-^{11}\},
\]
where we use interval upper bounds for the coefficient absolute
values. The recurrence \eqref{uhe:recurrence} gives
\[
 |c_j(a)|\le L_a a_-^{1-j},\qquad j\ge11.
\]
The induction step holds whenever \(j\ge a_-/4\), which is true
in this range for both choices of \(a\). For \(|h|\le\eta H\),
\(\eta\le\eta_0=1/40\) and \(q=\eta_0H/a_-<1\), summing the
tails gives
\[
 \begin{split}
 \left|f_a(h)-\sum_{j=0}^{12}c_j(a)h^j\right|
 &\le
 \eta^{13}L_a a_-\frac{(H/a_-)^{13}}{1-q},\\
 \left|f_a'(h)-\sum_{j=1}^{12}jc_j(a)h^{j-1}\right|
 &\le
 \eta^{12}L_a(H/a_-)^{12}
       \left(\frac{13}{1-q}+\frac q{(1-q)^2}\right).
 \end{split}
 \tag{UHE3}\label{uhe:tail}
\]
Since \(f_\alpha(h_x)\) and \(f_\zeta(h_y)\) are both \(O(\eta)\),
the tail in \(D(\eta,z)\) is \(O(\eta^{13})\). After dividing by
\(\eta^3\), the bounds in \eqref{uhe:tail} give an explicit
constant times \(\eta_0^{10}\).

We expand the truncated determinant in powers of \(\eta\), use
the leading quotient coefficient in \eqref{uhe:normalised-D}
and bound the higher coefficients on the closed disk
\[
 |z-z_c|\le r,\qquad
 z_c=-1.30655+2.26302\,i,\quad r=1/10,
\]
We then add the determinant tail from \eqref{uhe:tail}.
Take \(C_0\) and \(\alpha_0\) to be the midpoints of the brackets
for \(N(N-1)/12\) and \(\alpha\), respectively. On the boundary of
the disk, the reference cubic \(C_0z^3-\alpha_0^2\) has modulus
at least
\[
 3C_0|z_c|^2r-3C_0|z_c|r^2-C_0r^3
       -|C_0z_c^3-\alpha_0^2|.
\]
We use complex interval bounds in this expression. The total
parameter, polynomial and tail error is less than the contour
bound, with relative margin greater than \(0.5229\). Positivity
of the contour bound also allows us to compare the reference cubic
with \(3C_0z_c^2(z-z_c)\). Applying Rouch\'e's theorem twice gives
one root in the disk for every \(0\le\eta\le1/40\). For
\(\eta>0\), the corresponding \(x\) has positive imaginary part
and \(|\sqrt n\,x|<14\), so the disk count identifies the upper
branch.

Since \(x-\alpha=\eta z\), we bound the rescaled quantities
\[
 K=\eta\kappa
   =\frac{2x f_\alpha'(\eta z)}{f_\alpha(\eta z)/\eta},
 \qquad
 \widetilde Z=\frac{x}{K}
   =\frac{f_\alpha(\eta z)/\eta}{2f_\alpha'(\eta z)}.
\]
Both formulas extend to \(\eta=0\) in our disk. Using degree-\(12\)
series and \eqref{uhe:tail} gives the enclosures
\[
 \begin{array}{c|cc}
       &\operatorname{Re}&\operatorname{Im}\\ \hline
 K&[-2.724,-1.758]&[-4.554,-3.231]\\
 \widetilde Z&[-0.720,-0.591]&[1.061,1.196].
 \end{array}
\]
The rational bounds give margins greater than \(83.11\) and
\(3.6496\) in the two squared cone inequalities, and give
\(\operatorname{Re}(\widetilde Z^2)<-0.6079\).
Thus \(-\operatorname{Re}K>|K|/6\),
\(-\operatorname{Re}\widetilde Z>|\widetilde Z|/4\) and
\(\operatorname{Re}(\widetilde Z^2)<0\). Since \(\eta>0\), the
same cones contain \(\kappa=K/\eta\) and
\(x/\kappa=\eta\widetilde Z\).

It remains to check that the two contrast intervals overlap.
The Bessel brackets give
\[
 N-\frac1{64000}
 <\left(\frac{459083}{200000}\right)^2<N.
\]
All the interval comparisons above are verified in
\repofile{subresonant/unit_midrange/cap/hessian_endpoint_check.py}{endpoint checker}.
\end{proof}

\section{The radial calibration certificates}\label{sup:radial}

The comparison for one-signed minimisers above resonance comes
down to the signs of \(\bar u,p,\phi,T,a\) and \(4a-D^2\) for
\(0<\eta<1\) and \(0<r<1\). Some of these quantities vanish
at the endpoints. To obtain strict bounds on closed rectangles,
we first factor out the zeros.

\subsection{A regular parameter for the whole radial branch}

Set \(E(z)=J_0(\sqrt z)+\sqrt zJ_1(\sqrt z)\).
For \(\eta\in[0,1]\), let \(\lambda(\eta)\in[0,\alpha]\)
and \(\mu(\eta)\in[\beta,\zeta]\) be the two branches of the Robin
equation
\[
 (1-\eta)J_0(\sqrt z)-\eta\sqrt zJ_1(\sqrt z)=0 .
\]
For \(0<\eta<1\), the radial logarithmic derivative
\(\sqrt z J_0'(\sqrt z)/J_0(\sqrt z)\) takes the same value
\(1-1/\eta\) at these two roots. As \(\eta\) increases from \(0\)
to \(1\), the roots run from \(\alpha\) to \(0\) and from \(\zeta\)
to \(B=\beta\), respectively. Both branches are analytic at
\(\eta=0,1\). In the interior each branch \(x\) satisfies
\[
 x_\eta=-\frac{2x}{(1-\eta)^2+x\eta^2},
\]
and therefore
\[
 \frac{\dd}{\dd\eta}\log\frac{\mu}{\lambda}
 =\frac{2(\mu-\lambda)\eta^2}
 {[(1-\eta)^2+\lambda\eta^2][(1-\eta)^2+\mu\eta^2]}>0.
\]
Thus the ratio \(n=\mu/\lambda\) increases from \(N\) to \(\infty\),
with \(\lambda=\ell(n)\). The parameter \(\eta\) keeps both roots
bounded even as \(n\) tends to infinity. Let \(\kappa=1-1/\eta\)
be the common logarithmic derivative of the two Bessel factors.
The profile
\begin{equation}\label{eq:all-profile}
 \bar u(r)=\frac1{\mu-\lambda}
 \left\{\frac{J_0(\sqrt\lambda r)}{E(\lambda)}
       -\frac{J_0(\sqrt\mu r)}{E(\mu)}\right\}
\end{equation}
satisfies
\[
 \bar u(1)=\bar u'(1)=0,\qquad
 \bar u''(1)=\eta,\qquad \bar u'''(1)=-1,
\]
\[
 \Delta^2\bar u+c\Delta\bar u+d\bar u=0,
 \qquad c=\lambda+\mu,\quad d=\lambda\mu.
\]
For \(0<\eta<1\), the identity \(\eta/J_0(\sqrt x)=1/E(x)\)
on the Robin curve gives \(\bar u=\eta U\), where \(U\) is
the radial disk field normalised by \(U''(1)=1\).
The bounds from \path{robin_denominators.py},
\[
 E(\lambda)>\tfrac12,\qquad E(\mu)<-\tfrac14,
 \qquad \mu-\lambda>8,\qquad \sqrt\mu<6,
\]
give a joint analytic continuation of \(\bar u\) in \(r\) and
\(\eta\) to both endpoints. The check \path{robin_normalization_bounds.py}
gives \(\sqrt\alpha J_1(\sqrt\alpha)>1\).
Combining these bounds with the integral representation of \(J_0\),
we obtain, for every integer \(j\ge0\),
\begin{equation}\label{eq:all-uniform-derivatives}
 |\bar u^{(j)}(r)|\le6^j,\qquad
 |\phi^{(j)}(r)|\le3^j,
 \qquad
 \phi(r)=\frac{J_0(\sqrt\alpha r)}{\sqrt\alpha J_1(\sqrt\alpha)}.
\end{equation}
We use the estimate on the combined profile \(\bar u\) to bound
the Taylor tails. At the boundary we use two auxiliary solutions
to separate the dependence on \(\eta\) in the Taylor coefficients.

\subsection{The calibration and its exact identities}

The calibration uses the ratio \(t=\bar u/p\), where \(p\) is
the negative radial derivative of \(\bar u\). For \(0<r<1\) set
\begin{gather*}
 p=-\bar u',\quad A=-p'/p,\quad C=-1/r,\quad
 t=\bar u/p,\quad T=tA-1,\\
 H(r)=\int_0^r s\bar u(s)\dd s,\quad R=-dH/(rp),\\
 f=\eta\phi/p,\quad w=\eta\phi'/p,\quad
 z=(w+Af)/T,\quad g=f-tz,\quad a=1-g^2,\\
 G=r^2\{w^2-Cfw+\alpha f^2-A_r-(1+Ct)R+dt^2\},\\
 b=-tR-A-GC-fw,\quad D=2G-tG_r/T.
\end{gather*}
The signs established below give \(t_r<0\) and \(t\) runs from
infinity to zero for each fixed \(0<\eta<1\). This lets us regard
these quantities as functions of \(t\). The radial equations are
\begin{gather*}
 A_r=R+A^2+AC+c-C^2,\quad C_r=C^2,\quad t_r=T,\quad
 R_r=(A+C)R-dt,\\
 f_r=w+Af,\qquad w_r=(A+C)w-\alpha f.
\end{gather*}
Substituting these radial equations gives
\begin{align}
 G_r&=(2A-C)G+2C-A-tR-fw=b+2AG+2C,\nonumber\\
 b&=-G_t-DA-2C,\nonumber\\
 -2b+tb_t-2aA-DC&=-2zg,\label{eq:all-identities}\\
 aA^2+C^2+DAC-b_t&=-c+dt^2-z^2+\alpha f^2.\nonumber
\end{align}
To apply the square completion we need
\begin{equation}\label{eq:all-PSD}
 \bar u,p,\phi>0,\qquad T<0,\qquad
 a>0,\qquad 4a-D^2>0\qquad(0<r<1).
\end{equation}
The last two inequalities give positive definiteness of the
quadratic form in the completion.

\subsection{Interval bounds at the centre and boundary}

How do we check \eqref{eq:all-PSD} near the endpoints? At the
centre we use the evenness of the radial fields. Near \(r=1\) we
factor their zeros and resolve the joint limit \(\eta=0,r=1\).

In the centre chart set $s=r^2$, $\bar u=V(s)$ and $\phi=W(s)$.
Define
\begin{gather*}
 t_0=-\frac{V}{2V'},\qquad A_0=-1-\frac{2sV''}{V'},\qquad
 T_0=t_0A_0-s,\\
 f_0=-\frac{\eta W}{2V'},\qquad
 w=-\frac{\eta W'}{V'},\qquad
 R=\frac{d}{4V'}\frac1s\int_0^sV.
\end{gather*}
The nonsingular formulas are
\begin{align*}
 g^2&=\frac{s(f_0+t_0w)^2}{T_0^2},\\
 G&=1+A_0-A_0^2-cs-(2s-t_0)R+dt_0^2
             +sw^2+f_0w+\alpha f_0^2,\\
 D&=\frac{-G(2s+t_0)+t_0(t_0R+A_0+f_0w+2)}{T_0}.
\end{align*}
We evaluate these expressions on \(0\le\eta\le1\), \(0\le s\le9/16\)
using Bessel series and interval bounds on their derivatives. Define
\[
 \mathcal B_m(z)=\sum_{k=0}^{\infty}\frac{(-z/4)^k}{k!(k+m)!},
 \qquad \mathcal B_m'=-\mathcal B_{m+1}/4.
\]
Every Bessel argument lies in $[0,36]$. The degree-$K$ truncation,
with $K=30$, has absolute tail at most
\[
 \mathcal T_K=\frac{9^{K+1}}{((K+1)!)^2[1-9/(K+2)^2]}.
\]
Poisson's integral gives \(|\mathcal B_m(z)|\le1/m!\) for real \(z\ge0\).
For an interval argument of width \(w_*\), we evaluate the
truncation at the midpoint and widen the enclosure by
\(w_*/[8(m+1)!]+\mathcal T_K\). The derivative bounds follow by
applying the same estimates to \(\mathcal B_m'\) and
\(\mathcal B_m''=\mathcal B_{m+2}/16\).

To bound \(a\) and \(4a-D^2\), we first enclose
\(h=(f_0+t_0w)/T_0\) and \(D\), using
\[
 h=\frac{\eta(VW'-WV')}{VV'+2sVV''-2s(V')^2},\qquad a=1-sh^2.
\]
For a rectangle \(Q\) with centre \(q_0\) and coordinate
increments \((e,z)\), Taylor's formula with integral remainder
encloses each of \(L=h,D,-T_0,V,-V'\) in the interval
\[
 L(q_0)+L_\eta(q_0)e+L_s(q_0)z
 +\tfrac12L_{\eta\eta}(Q)e^2+L_{\eta s}(Q)ez
 +\tfrac12L_{ss}(Q)z^2.
\]
We differentiate in exact interval arithmetic and check the
intermediate denominators before dividing. This gives
\(a>87/1000\) and \(4a-D^2>3/10000\) on the centre chart.

At the boundary we use \(s=1-r\) and factor the fields as
\[
 \bar u=s^2U,\quad p=sP,\quad\phi=sF,\quad
 Q=P+sP_s,\quad L=Q_s,\quad J=F+sF_s,
 \quad Z=P^2-UQ.
\]
Integration of the radial equation gives
$dH=1-\eta-ds^3K$ with $K$ analytic. Then
\[
 T=-Z/P^2,\qquad g=\eta(FP-UJ)/Z,\qquad G=s\mathcal W/P^2,
\]
where
\begin{align*}
 s^3\mathcal W={}&r^2(\eta^2J^2-Q^2)
  +s(r^2PL-r\eta^2FJ)+(1-\eta)s(rP-sU)\\
 &+\alpha\eta^2s^2r^2F^2
  +ds^4\{r^2U^2-K(rP-sU)\}.
\end{align*}
The first three coefficients of \(s^3\mathcal W\) vanish
exactly, and \(\mathcal W(\eta,0)=\eta(1+\eta)/3\).
We compute the remaining coefficients from
\[
 (1-s)s\mathcal W_s+(3-2s)\mathcal W
 =(1-s)\frac{QP-\eta^2FJ}{s}+2P^2-(1-\eta-ds^3K)U.
\]
At fixed $c,d$, the Taylor coefficients of $\bar u$ are affine
in $\eta$, and those of $\mathcal W$ are quadratic in $\eta$.
Separating the powers of \(\eta\) reduces interval overestimation.

Near \(\eta=s=0\) we distinguish \(s\le\eta\) and
\(\eta\le s\). In the first region we set \(s=\eta x\) and use
\(0\le\eta\le1/4\), \(0\le x\le1\). Factoring
\(U=\eta\widetilde U\), \(P=\eta\widetilde P\),
\(Q=\eta\widetilde Q\), \(\mathcal W=\eta\widetilde{\mathcal W}\)
leaves analytic fields. We check positivity of \(a/x^2\),
\((4a-D^2)/x^2\) and
\(\widetilde Z=\widetilde P^2-\widetilde U\widetilde Q\).
At \(\eta=0\), let \(q=x^2+4x+6\). Then
\[
 g=\frac{2(2x+3)}q,\quad
 a=\frac{x^2(x+2)(x+6)}{q^2},\quad
 G=\frac{x(x+4)}{3(x+2)^2},\quad
 D=\frac{2x(x^2+6x+6)}{3(x+2)q},
\]
\[
 \frac{4a-D^2}{x^2}
 =\frac{4(8x^4+96x^3+384x^2+648x+396)}{9(x+2)^2q^2}>0.
\]

In the second region we set \(\eta=s\delta\) and use
\(0\le s\le1/4\), \(0\le\delta\le1\). We factor
\(U=s\widehat U\), \(P=s\widehat P\),
\(Q=s\widehat Q\), \(K=s\widehat K\),
\(\mathcal W=s\widehat{\mathcal W}\). With
$\widehat Z=\widehat P^2-\widehat U\widehat Q$ and
$M_0=1-s\delta-ds^4\widehat K$, one has
\begin{align*}
 g&=\delta(F\widehat P-\widehat UJ)/\widehat Z,\\
 D&=\frac{1}{\widehat Z}\left\{
 2\widehat{\mathcal W}
 +\frac{M_0\widehat U^2}{r\widehat P}
 +\frac{\widehat U(\delta^2FJ-\widehat Q\widehat P)}{\widehat P}
 +\frac{s\widehat U\widehat{\mathcal W}}{r\widehat P}
 -\frac{2s\widehat U\widehat P}{r}\right\}.
\end{align*}
Here \(\widehat{\mathcal W}(0,\delta)=(1+4\delta)/12\).
We can check \(a\) and \(4a-D^2\) directly.
Away from the corner we use \(1/16\le\eta\le1\) and
\(0\le s\le1/4\), splitting the interval for \(\eta\) at \(1/4\).
In both parts we use
\[
 \frac Ds=\frac{2\mathcal W}{P^2}
 -\frac{UP}{Z}\left\{\frac{\mathcal W}{P^2}
   +\frac{s\mathcal W_s}{P^2}-\frac{2s\mathcal WP_s}{P^3}\right\}.
\]
The numerator \(Z^2-\eta^2(FP-UJ)^2\) has a double zero at
\(s=0\). After factoring it out we check that \(a/s^2\),
\((4a-D^2)/s^2\) and \(Z\) are positive. These two corner
charts and the regular region cover the boundary rectangle.

It remains to bound the Taylor errors after the divisions.
Let \(a_*=1/4\) and define
\[
 E_C(M)=\frac{432\,6^{M+1}}
 {(M+1)!\,[1-6a_* /(M+2)]}.
\]
The raw degree-$M$ field tails are at most $E_C(M)s^{M+1}$,
and their first derivatives at most $2(M+1)E_C(M)s^M$.
In the chart \(s=\eta x\) we use \(M=48\) and allow the common error
\[
 \mathcal E_A=10^{30}E_C(48)a_*^{42}<5.1\times10^{-18}.
\]
We bound the fields by \(10^6\) and expand the difference in each
quadratic or quartic product one factor at a time. After division
by the exact vanishing factors and the required differentiation,
this bounds the errors by \(10^{22}E_C(48)a_*^{42}\), which is
strictly below \(\mathcal E_A\).

For the other boundary charts we use \(M=38\) and
\[
 \mathcal E_E=10^{22}E_C(38)a_*^{32}<2.658\times10^{-11}.
\]
The recurrence gives \(\mathcal W\) through degree \(M-3\).
To bound the tail, we majorise the absolute coefficients of
\(s^3\mathcal W\) by those of
$P_6(s)e^{6s}+P_{12}(s)e^{12s}$, where
\begin{align*}
 P_6(s)={}&9(1+s)^2+45s(1+s)+18s^2+54s^2(1+s)^2,\\
 P_{12}(s)={}&1296(1+s)^2+7776s(1+s)^2
       +64800s^4(1+s)^2\\
       &+57600s^4(1+s)+28800s^5.
\end{align*}
These bounds follow from
$U\preceq18e^{6s}$, $P,Q\preceq36e^{6s}$,
$L\preceq216e^{6s}$, $F,J\preceq3e^{3s}$,
and $K\preceq8e^{6s}$, using the displayed numerator for
$\mathcal W$. Here $\preceq$ denotes coefficientwise domination
of absolute values. Define
\[
 T(z,k)=\frac{z^k}{k!\,[1-z/(k+1)]},\qquad K_*=M+1,
\]
and set
\begin{align*}
 N_v&=\sum_{q\in\{6,12\}}\sum_j
        P_{q,j}a_*^jT(qa_*,K_*-j),\\
 N_d&=\sum_{q\in\{6,12\}}\sum_jP_{q,j}a_*^j
    \{qa_*T(qa_*,K_*-j-1)+jT(qa_*,K_*-j)\}.
\end{align*}
The omitted $\mathcal W$ and derivative terms, including the
resolved chart divisions, are bounded by
\[
 4N_v/a_*^3+N_d/a_*^4<2.748\times10^{-16}<\mathcal E_E/100.
\]
The other raw-field and polynomial errors are bounded by
$10^{20}E_C(38)a_*^{32}=\mathcal E_E/100$.
In the two regular charts we remove the double zero from
\(Z-\eta(FP-UJ)\) before multiplying by \(Z+\eta(FP-UJ)\).
This lets us perform the cancellation in a quadratic expression.
The common error bound covers the truncations in both charts,
labelled E and Esmall in the certificate index.

In the chart \(\eta=s\delta\), labelled B in the certificate
index, \path{all_contrast_endpoint_B_covered.py} calls
\texttt{chart\_B} in \path{all_contrast_endpoint.py} with the
degree-$38$ coefficients from \path{all_contrast_endpoint_fast.py}.
We enclose \(\widehat U,\widehat P,\widehat Q,F,J,\widehat K\)
and \(\widehat{\mathcal W}\) separately, checking
\(\sum_j|p_j|a_*^j<10^6\) for each polynomial, where \(|p_j|\)
is the upper bound on the absolute coefficient. We widen each
interval by \(\mathcal E_E\) at both ends. For the first six fields,
division by at most one power of \(s\) leaves a tail of at most
$E_C(38)a_*^{38}<6.489\times10^{-37}$. For
$\widehat{\mathcal W}=\mathcal W/s$, the tail is at most
$N_v/a_*^4=4N_v/a_*^3<\mathcal E_E/100$.
From these intervals we form \(\widehat Z\) and check that
\(\widehat U,\widehat P,F,\widehat Z\) are positive.
Since \(r\ge3/4\), we can then divide in the formulas for \(g\)
and \(D\). A rectangle passes when \(a=1-g^2\) and \(4a-D^2\)
also have positive lower bounds.

The degrees in the table refer to the Bessel series at the centre and
to the raw fields at the boundary.
\begin{center}
\begin{tabular}{lrr}
\hline
Chart & Degree & Certified rectangles\\
\hline
$0\le r^2\le9/16$, $0\le\eta\le1$ & 30 & 27167\\
$s=\eta x$, $0\le\eta\le1/4$, $0\le x\le1$ & 48 & 433\\
$1/4\le\eta\le1$, $0\le s\le1/4$ & 38 & 1574\\
$1/16\le\eta\le1/4$, $0\le s\le1/4$ & 38 & 773\\
$\eta=s\delta\le1/16$ in the resolved chart & 38 & 2073\\
\hline
\end{tabular}
\end{center}
Of the 2073 rectangles in the last row, 94 cross
\(s\delta=1/16\). On \([s_-,s_+]\times[\delta_-,\delta_+]\) we
enclose the Robin roots for
\(s_-\delta_-\le\eta\le\min\{s_+\delta_+,1/16\}\)
and evaluate the fields on the full \(s\) and \(\delta\) intervals.
The regular charts cover the part with \(s\delta\ge1/16\) and all
68 remaining rectangles. All sign tests use directed integer arithmetic.
The \repofile{certificates/continuation/audit_all_contrast_coverage.py}{coverage calculation}
uses the rational endpoints of the five partitions to check that
their union covers \((\eta,r)\in[0,1]^2\), including the common
edge \(r=3/4\) of the centre and boundary charts.


\end{document}
